\section{Discussion: two treatments of a sparse \Seg{} token}
\label{sec:discussion}

The empirical core is that a 20-D corridor can be rasterized, and that the sparse \Seg{} prompt DeSa2VA locates \emph{inside} the token can be addressed from the \emph{decoder} side, not only by splitting the token.

\paragraph{Carrier replacement vs.\ DeSa2VA's architectural decoupling.}
DeSa2VA \citep{dang2025desa2va} splits a Sa2VA hidden state into text and visual streams \emph{before} SAM2, because a single \Seg{} token cannot carry both semantics and a dense visual prompt.
ECF never splits that state and never asks SAM2 to become a better prompt reader; it lets the entangled token bind a caption to an object---what language models already do---and writes the dense spatial structure as $\phih$.
Two independent observations support this reading.
First, A1: the same $h$ that rasterizes a rail is a useless SAM2 prompt (2.43 gIoU), so the information is in the ParamHead$+\cR$ path, not in some latent SAM2-compatible embedding.
Second, DeSa2VA's own ablation shows that extra text-derived mask supervision into an entangled \Seg{} pathway hurts Sa2VA (RefCOCO 78.9$\to$73.4 without decoupling).
We do not feed extra text into SAM2; we stop feeding $h$ into SAM2 on the evaluated path, and we are not claiming to reproduce DeSa2VA's RefCOCO gains, only that the sparse-prompt diagnosis has a decoder-side, geometric treatment as well as an encoder-side one.

\paragraph{Low-dimensional corridors vs.\ free masks.}
Parametric lane heads \citep{tabelini2021polylanenet,feng2022bezierlanenet} already treat ego-view rails as polylines.
What is new here is the \emph{binding}: the polyline is the meaning of a referring \Seg{} token inside an MLLM, trained with a raster closed loop so that language gradients and geometry gradients share $h$.
Figure~\ref{fig:framework} draws this forward path (Section~\ref{sec:method}), and Figure~\ref{fig:display} shows the shapes the field must cover across this split.
A2 shows the loop is load-bearing: without $\mathcal{L}_{\mathrm{mask}}$, \PrAt{0.7} drops to 54.81 even though a parameter loss is present.

\paragraph{Carrier versus backbone.}
The control that isolates the carrier is internal to our system rather than external: A1 reuses the same trained LoRA and ParamHead and changes only where the mask comes from, giving 2.43 versus 80.67 gIoU.
That already undercuts the trivial reading that the masks come from the language backbone.
The carrier claim itself stands on A1 and on the failure mode of Figure~\ref{fig:qual} row 4, where high-IoU masks land on the wrong rail: on this target, decoder strength by itself did not deliver target selection; changing the carrier did.

\paragraph{What this paper does and does not replace.}
The two treatments are not substitutes: decoupling improves the encoder-side prompt that a generalist SAM2 receives, while ECF assumes a corridor target and changes the decoder.
We read them as alternatives for different target classes rather than as a refutation of either.
